\documentclass[11pt,a4paper]{article}
\usepackage[margin=0.82in]{geometry}
\usepackage[T1]{fontenc}
\usepackage{lmodern,microtype,amsmath,amssymb}
\usepackage[hidelinks]{hyperref}
\usepackage{enumitem}
\setlength{\parindent}{0pt}
\setlength{\parskip}{5pt}
\setlength{\emergencystretch}{2em}
\setlist[itemize]{leftmargin=1.2em,itemsep=3pt,topsep=3pt}
\usepackage{titlesec}
\titleformat{\section}{\large\bfseries}{\thesection.}{.5em}{}
\titlespacing*{\section}{0pt}{10pt}{5pt}
\pagestyle{plain}
\begin{document}
\begin{center}
{\LARGE\bfseries Understanding the proof of Barnette's conjecture}\par
\vspace{4pt}
{\large Mathathon, Combinatorics}\par
\vspace{5pt}
Convex Curiosity\par
Anouk Brose (UC Davis), Tom de Groot (independent), Santiago Morales Duarte (UC Davis)\par
\vspace{4pt}
{\footnotesize
\url{https://github.com/Tomodovodoo/mathathon-Barnette-s-Conjecture}\par
Code snapshot: \texttt{c39082963fea936f98b3aeb7f6d1f9258ec82c2e}\par
15 November 2026}
\end{center}
\vspace{-2pt}

\textbf{Abstract.} We set out to make the proof of Barnette's conjecture digestible through interactive tools, a simpler construction and working code. We built an interactive website and proof walkthrough, shortened the cycle-recovery route from six steps to four, and implemented algorithms that find the cycle. We tested simpler routes and studied which assumptions each step needs. This report describes those contributions, our empirical findings, and how we used AI. The exposition supports the tools and explains their mathematics.

\section{Our goal}
We wanted to understand why the proof of Barnette's conjecture works. The theorem says that every cubic, bipartite, $3$-connected planar graph has a Hamiltonian cycle; equivalently, every simple $3$-polytope with even-sided faces has a Hamiltonian graph~\cite{BarnetteConj}. Tait hoped to prove the four-colour theorem through Hamiltonian cycles in simple $3$-polytopes~\cite{Tait}, Tutte found a non-Hamiltonian one~\cite{Tutte}, and Barnette added the bipartite assumption. Goodey proved the case with only quadrilateral and hexagonal faces~\cite{Goodey}, and Gorsky, Steiner and Wiederrecht related the conjecture to a strengthening with Pfaffian graphs~\cite{GSW}. We read OpenAI's proof~\cite{OpenAI}, Li's account~\cite{Li}, and a collaborator's constructive notes~\cite{Constructive}. We also looked at online explanations, including automatically generated digestions.

We found steps that needed a clearer explanation, a more direct route, or a link to a standard graph-theoretic tool. We wanted readers to distinguish the main ideas from the details that support them. We also wanted to make each step's assumptions explicit, especially where an assumption determines which route can work.

Our background in polytopes and polyhedral graphs shaped the project. We wanted to follow each change of viewpoint from the graph to the matchings, trees and counts that solve the problem. Understanding these objects takes work. We aimed to explain what each object does before asking the reader to use it.

We studied the assumptions so that others could reuse and generalise the lemmas. A proof may state a lemma under all the assumptions of the main theorem even when the lemma needs fewer. We tried to separate those needs. We also looked for known theorems that could explain a step or replace it with a simpler argument.

Our main contributions are the interactive website, the shorter cycle construction, and the implementation. The website lets readers explore the cycle on a polytope. The proof walkthrough shows how the construction works. The code turns the proof into a procedure that finds a cycle and lets us test alternative choices. Before building them, we checked most of the steps of the preprint against the other sources and our computations. We developed these together during the event. The exposition records the supporting arguments and connections to earlier work. Our exposition has not undergone journal peer review.

\newpage
\section{What we simplified and clarified}
We shortened the route from a good state to a Hamiltonian cycle from six steps to four. This count refers to the cycle-recovery construction. The existence argument supplies its starting state, and the reduction along separating triangles handles the remaining cases.

The earlier six-step presentation built a tree of faces before selecting edges. It then used the disk identity to prove that the selection had no cycle, recovered the Hamiltonian cycle, and deduced the two induced trees. We replace the intermediate tree of faces with a colouring and use an innermost-cycle argument to prove acyclicity. This removes both that intermediate construction and the disk identity from recovery. We connect the result to Li's directed-trace argument and Stein's induced-tree characterisation~\cite{Li,Stein}.

Let $G$ be the original graph and $T$ its dual triangulation. We fix an outer black face $t_0$ and call its corners roots. A Tutte state $s$ chooses a corner in each other black face, using every nonroot vertex once. The sides opposite those corners form $Q_s$. A good state makes $Q_s$ a forest with one root in each component.

\begin{enumerate}[leftmargin=1.5em,itemsep=3pt,topsep=3pt]
\item To avoid a prescribed edge $e$ incident to the vertex dual to $t_0$, add the two sides of $t_0$ other than the side dual to $e$. Together with $Q_s$ they form a spanning tree of $T$.
\item Give adjacent vertices of this tree opposite colours. Select all edges of $T$ whose endpoints share a colour.
\item Use the innermost-cycle argument to show that the selected edges form two induced trees, one for each colour.
\item Take the dual edges of the unselected edges. They form the Hamiltonian cycle in $G$.
\end{enumerate}

We separated this construction from the proof that a good state exists. A max-flow argument supplies a pair of disjoint states $(r,s)$ when $T$ has no separating triangle~\cite{OpenAI}. The partner matters: with a partner, acyclicity of $Q_s$ ensures that every component has one root. Without it, a forest alone need not give a good state.

The existence proof groups one weighted sum in two ways. Fixing $s$ cancels terms whenever $Q_s$ has a cycle. The directed graph $D(r,s)$ has an edge from $r(t)$ to $s(t)$ for each black face $t\ne t_0$. Grouping by its cycles gives one factor $1-z^{a}$ per cycle, with $a>0$. The contributions of lowest order at $z=1$ come from the $2$-factors with the fewest cycles and share a sign. A good state must remain. The constructive notes already give the polynomial method~\cite{Constructive}; we explain its geometry through crossings and examples.

When $T$ has a separating triangle, we cut it into smaller pieces, apply induction and glue the cycles. We keep the stronger form that lets us prescribe an edge for a cycle to avoid, so the pieces fit. This distinguishes the constructive reductions from the step that proves existence.

\section{What we tried and learned}
We tried to replace the signed count with tree and matching theorems. Tutte-state lattices~\cite{HineKalman} suggest choosing extreme states. The theorems of Rado, Nash-Williams and Edmonds suggest choosing or packing trees. A state must choose one corner in each black face other than $t_0$ and use each nonroot vertex once. Our tree-packing formulation gives spanning trees that meet either condition separately within a fixed $2$-factor of the face--corner incidence graph. A single tree need not meet both.

We used counterexamples to find where the shortcuts fail. An $8$-vertex dual example shows that an acyclic opposite-edge graph can put two roots in one component and none in another. A $16$-vertex example has a $2$-factor whose two orientations both fail the tree condition. Other examples defeat strict descent and unconditional lifting after contraction. These tests left sharper questions about moves that preserve the cycle count and contractions that preserve enough structure. We identified the missing steps, but did not complete a new proof without signs.

\newpage
\section{What we built and tested}
We built an interactive website to make the cycle a visible object that readers can explore. In \emph{Every Vertex Once}, a runner traces a Hamiltonian cycle on a rotating polytope. This gives a direct view of the original problem. The separate proof walkthrough follows the mathematical construction and the objects it introduces. Together, the two tools connect the cycle we want with the steps that produce it.

We built and revised the visual explanations alongside the mathematics. The cube follows the construction from start to finish. A shared example shows the two cycle systems, while a crossing diagram explains the weight change. The route map distinguishes successful, failed and open approaches. We revised layouts, labels and scales as the explanation changed. Scripts check that the states and cycles in the drawings match the mathematics.

We implemented two algorithms based on the proof that return a Hamiltonian cycle. The search chooses corners one face at a time. It rejects a partial choice when no matching can complete it or when the opposite edges already violate the rooted-forest condition. A second procedure follows the proof's exchange moves. These implementations build on the constructive notes~\cite{Constructive}. They let us test choices that the existence proof leaves open and show the construction at work.

We tested the code against independent implementations and direct checks. We compared Python and C versions, enumerated small cases, and tested constructed families and random graphs. We checked that each returned cycle visits every vertex exactly once and avoids the prescribed edge. We also used counterexamples to show when a proposed shortcut lost a necessary condition.

We measured search cost using the number of retained partial choices divided by $n=|V(G)|/2-1$. On sampled random graphs, the fitted median grew by a factor of $1.056$--$1.064$ per unit increase in $n$ with breadth-first face order and $1.15$--$1.22$ with the stored order. The proven worst-case upper bound is $(\sqrt{3})^n$ times a polynomial factor. These fits describe the sampled medians. The supplementary experiment notes give the ranges and cost measure.

We used the supporting exposition to connect the tools, the simplification and the tests. We compared the source arguments, mapped their dependencies, placed definitions beside their first use, and separated construction from existence. We studied the assumptions so that others can generalise the lemmas. The failed tree-packing and contraction routes identify conditions that further work must address. The exposition records this reasoning; the website, walkthrough and code give readers other ways to understand and use it.

\newpage
\section{AI use, attribution, and outcome}
We used Claude and GPT models to help understand the proof, explore simpler routes, and explain the mathematics through text and visualisations. In this project, we used them to compare arguments, connect steps to known theorems, develop examples, write code and review explanations. We checked their suggestions against the source papers, mathematical arguments and computations.

\begin{itemize}
\item We asked AI to explore alternative constructions, weaker assumptions and links to known theorems. This helped develop the direct recovery argument and proposals for descent and tree-based proofs. Counterexamples ruled out several proposals.
\item We used AI to compare the source accounts and review drafts. This helped us separate construction from existence, explain the two groupings of the count, and introduce definitions when readers need them.
\item We used AI to help write code, draw examples and run experiments. The examples show how weights change and terms cancel. Automated checks test that the pictures match the mathematics.
\item We asked for separate mathematical reviews and used independent computations and Lean checks. These checks exposed missing assumptions, overstrong claims and attribution errors. We corrected them before using the results.
\end{itemize}

We tried Open Math OS with open models to seek further proof continuations~\cite{MathOS}. Those runs produced incorrect or unsupported claims. We accepted no useful continuation from them.

Several questions remain open for us. Proposition~3.1, Lemma~4.1 and Lemma~6.1 of the preprint~\cite{OpenAI} are stated under more assumptions than their proofs appear to use: the disk identity uses only a state and the triangulated disk, and the root count uses only the outdegrees that the partner supplies. We want to drop the unneeded assumptions and give examples, such as our $8$-vertex example for Lemma~6.1, showing why the remaining ones are needed. Lemma~2.1 can be strengthened to an equivalence: a set of edges is a spanning tree exactly when the duals of the remaining edges form a spanning tree of the dual. We would like to contribute this tree--cotree duality and other reusable statements to mathlib, and to distinguish the standard tools from the new connections. Finally, our search is exponential in the worst case; whether a polynomial-time algorithm can find the cycle remains open.

We set the goals, chose which questions to pursue, made the mathematical and editorial decisions, and developed the interactive walkthrough. We used reviews to restore necessary assumptions, credit earlier arguments, narrow claims and correct walk counts that depend on the drawing. These internal checks do not amount to external peer review.

We credit the theorem and its central counting method to the earlier sources. Our work explains those ideas, shortens specific steps and tests proposed alternatives.

{\footnotesize
\begin{thebibliography}{99}
\setlength{\itemsep}{0pt}\setlength{\parskip}{0pt}
\bibitem{OpenAI} OpenAI, \emph{Paired states and Hamiltonian cycles in cubic bipartite planar graphs}, preprint, 24 September 2026.
\bibitem{Li} Z.~Li, \emph{A note on the proof of Barnette's conjecture}, arXiv:2610.30490, 2026.
\bibitem{Constructive} \emph{A constructive proof of Barnette's conjecture}, collaborator's working notes, 18 October 2026, supplied without an author line as \texttt{papers/barnette\_trees.pdf}.
\bibitem{Stein} S.~K.~Stein, \emph{B-sets and planar maps}, Pacific J. Math. \textbf{37} (1971), 217--224.
\bibitem{HineKalman} C.~Hine and T.~K\'alm\'an, \emph{Clock theorems for triangulated surfaces}, arXiv:1808.06091, 2018.
\bibitem{BarnetteConj} B.~Gr\"unbaum, Unsolved Problem 5 (conjecture of D.~W.~Barnette), in W.~T.~Tutte (ed.), \emph{Recent Progress in Combinatorics}, Academic Press, 1969, p.~343.
\bibitem{Tait} P.~G.~Tait, \emph{Listing's Topologie}, Philos. Mag. (5) \textbf{17} (1884), 30--46.
\bibitem{Tutte} W.~T.~Tutte, \emph{On Hamiltonian circuits}, J. London Math. Soc. \textbf{21} (1946), 98--101.
\bibitem{Goodey} P.~R.~Goodey, \emph{Hamiltonian circuits in polytopes with even sided faces}, Israel J. Math. \textbf{22} (1975), 52--56.
\bibitem{GSW} M.~Gorsky, R.~Steiner and S.~Wiederrecht, \emph{Matching theory and Barnette's conjecture}, Discrete Math. \textbf{346} (2023), 113249.
\bibitem{MathOS} Open Math Model, \emph{Open Math OS}, \url{https://www.openmathmodel.org/}, accessed 15 November 2026.
\end{thebibliography}}
\end{document}
